% GENERATED FILE. Edit canonical lessons, then run npm run generate:books.
% canonical-source-hash: fnv1a64:30f404881cc28bfc
% canonical-lessons: ZH-C531-tuoxian, ZH-C531-jihu, ZH-C531-jiaoxing, ZH-C531-buxing, ZH-C531-shishou, ZH-R531-review-526-531

\chapter{A Lucky Escape}
\label{ch:zh-a-lucky-escape}
\hlchaptermodality{531}

\begin{chapteropening}
\textbf{By the end of this chapter:} \emph{I can say it is a good thing that, all but, by a fluke and unfortunately, and talk about losing.}

Complete the last lesson of chapter 531: I can say it is a good thing that, all but, by a fluke and unfortunately, and talk about losing.
\end{chapteropening}

\section[tuōxiǎn]{tuōxiǎn — to get out of danger}

\label{lesson:ZH-C531-tuoxian}



\begin{quote}
Before the new one: say the Mandarin for a prize you win, then the Mandarin for the jackpot.
\end{quote}

\subsection*{You'll want to know: tuōxiǎn}
\emph{tuōxiǎn} — \textquotedblleft{}to get out of danger\textquotedblright{}.

\begin{quote}
\emph{zhǐyào zhǎo dào lù wǒmen jiù néng tuōxiǎn} — As long as we find the way, we'll get out of danger.
\end{quote}

Its characters come later; for now it is a word for the ear.

\begin{cousinweb}
\emph{tuōxiǎn} is built from \emph{tuō} and \emph{xiǎn} 'danger', both inherited from Middle Chinese, \emph{thwat} and \emph{xjemX} (Guangyun, in Baxter's transcription). The dictionary glosses the first as bone with the flesh taken off, and the second 'dangerous' and 'hard'; the sense 'to slip free' is not given for the first.
\end{cousinweb}

\begin{grammarlens}[title={What you've built}]
One more word for saying what would happen.
\end{grammarlens}

\subsection*{Your turn}
\begin{itemize}
\raggedright
  \item \emph{Say it:} \emph{tuōxiǎn}
  \item \emph{Say it:} \emph{tuōxiǎn}, once more
  \item \emph{Say it:} \emph{zhǐyào zhǎo dào lù wǒmen jiù néng tuōxiǎn} — As long as we find the way, we'll get out of danger.
  \item \emph{Recall:} say the Mandarin for a prize you win, then the Mandarin for the jackpot, then say \emph{tuōxiǎn} again
\end{itemize}

\begin{tcolorbox}[breakable,colback=teal!4,colframe=teal!35!black,title={Before you move on}]
What does tuōxiǎn mean? (\textquotedblleft{}To get out of danger\textquotedblright{}.) Say it once more.
\end{tcolorbox}
\section[jīhū]{jīhū — very nearly; all but}

\label{lesson:ZH-C531-jihu}



\begin{quote}
Before the new one: say the Mandarin for the jackpot, then the Mandarin for to get out of danger.
\end{quote}

\subsection*{You'll want to know: jīhū}
\emph{jīhū} — \textquotedblleft{}very nearly; all but\textquotedblright{}.

\begin{quote}
\emph{rúguǒ xiàyǔ gōngyuán lǐ jīhū méiyǒu rén} — If it rains, there's hardly anyone in the park.
\end{quote}

Its characters come later; for now it is a word for the ear.

\begin{cousinweb}
\emph{jīhū} is built from \emph{jī} and \emph{hū}, both inherited from Middle Chinese, \emph{kj+j} and \emph{hu} (Guangyun, in Baxter's transcription). The dictionary lists the first four times: \emph{kj+j}, which gives jī, in a phrase meaning 'very nearly', and \emph{gj+j}, glossed 'near'. It glosses the second as a grammar word; the level tone of hū is not the regular outcome of \emph{hu}, which would give hú.
\end{cousinweb}

\begin{grammarlens}[title={What you've built}]
One more word for saying what would happen.
\end{grammarlens}

\subsection*{Your turn}
\begin{itemize}
\raggedright
  \item \emph{Say it:} \emph{jīhū}
  \item \emph{Say it:} \emph{jīhū}, once more
  \item \emph{Say it:} \emph{rúguǒ xiàyǔ gōngyuán lǐ jīhū méiyǒu rén} — If it rains, there's hardly anyone in the park.
  \item \emph{Recall:} say the Mandarin for the jackpot, then the Mandarin for to get out of danger, then say \emph{jīhū} again
\end{itemize}

\begin{tcolorbox}[breakable,colback=teal!4,colframe=teal!35!black,title={Before you move on}]
What does jīhū mean? (\textquotedblleft{}Very nearly; all but\textquotedblright{}.) Say it once more.
\end{tcolorbox}
\section[jiǎoxìng]{jiǎoxìng — lucky by a fluke}

\label{lesson:ZH-C531-jiaoxing}



\begin{quote}
Before the new one: say the Mandarin for to get out of danger, then the Mandarin for all but.
\end{quote}

\subsection*{You'll want to know: jiǎoxìng}
\emph{jiǎoxìng} — \textquotedblleft{}lucky by a fluke\textquotedblright{}.

\begin{quote}
\emph{wǒmen zhè cì méi chūshì zhǐ shì jiǎoxìng} — Nothing happened to us this time, but it was only a fluke.
\end{quote}

Its characters come later; for now it is a word for the ear.

\begin{cousinweb}
\emph{jiǎoxìng} is built from \emph{jiǎo} and \emph{xìng}, both inherited from Middle Chinese, \emph{kewX} and \emph{heangX} (Guangyun, in Baxter's transcription). The dictionary reads the first character as written today \emph{ngew}, the name of a far country, and lists the word itself, in an older spelling, under its second syllable; the first syllable of that spelling, \emph{kewX}, regularly gives jiǎo.
\end{cousinweb}

\begin{grammarlens}[title={What you've built}]
One more word for saying what would happen.
\end{grammarlens}

\subsection*{Your turn}
\begin{itemize}
\raggedright
  \item \emph{Say it:} \emph{jiǎoxìng}
  \item \emph{Say it:} \emph{jiǎoxìng}, once more
  \item \emph{Say it:} \emph{wǒmen zhè cì méi chūshì zhǐ shì jiǎoxìng} — Nothing happened to us this time, but it was only a fluke.
  \item \emph{Recall:} say the Mandarin for to get out of danger, then the Mandarin for all but, then say \emph{jiǎoxìng} again
\end{itemize}

\begin{tcolorbox}[breakable,colback=teal!4,colframe=teal!35!black,title={Before you move on}]
What does jiǎoxìng mean? (\textquotedblleft{}Lucky by a fluke\textquotedblright{}.) Say it once more.
\end{tcolorbox}
\section[búxìng]{búxìng — misfortune; unfortunate}

\label{lesson:ZH-C531-buxing}



\begin{quote}
Before the new one: say the Mandarin for all but, then the Mandarin for lucky by a fluke.
\end{quote}

\subsection*{You'll want to know: búxìng}
\emph{búxìng} — \textquotedblleft{}misfortune; unfortunate\textquotedblright{}.

\begin{quote}
\emph{jíshǐ yùdào búxìng yě búyào fàngqì} — Even if misfortune comes your way, don't give up.
\end{quote}

Its characters come later; for now it is a word for the ear.

Tone: \zh{不} is said bú before the falling tone of xìng, and the spelling shows it.

\begin{cousinweb}
\emph{búxìng} is built from \zh{不} 'not' and \emph{xìng} 'fortunate', both inherited from Middle Chinese, \emph{pjut} and \emph{heangX} (Guangyun, in Baxter's transcription), which glosses the second as having good fortune and escaping harm.
\end{cousinweb}

\begin{grammarlens}[title={What you've built}]
One more word for saying what would happen.
\end{grammarlens}

\subsection*{Your turn}
\begin{itemize}
\raggedright
  \item \emph{Say it:} \emph{búxìng}
  \item \emph{Say it:} \emph{búxìng}, once more
  \item \emph{Say it:} \emph{jíshǐ yùdào búxìng yě búyào fàngqì} — Even if misfortune comes your way, don't give up.
  \item \emph{Recall:} say the Mandarin for all but, then the Mandarin for lucky by a fluke, then say \emph{búxìng} again
\end{itemize}

\begin{tcolorbox}[breakable,colback=teal!4,colframe=teal!35!black,title={Before you move on}]
What does búxìng mean? (\textquotedblleft{}Misfortune; unfortunate\textquotedblright{}.) Say it once more.
\end{tcolorbox}
\section[shīshǒu]{shīshǒu — to lose one's grip; to do by a slip of the hand}

\label{lesson:ZH-C531-shishou}



\begin{quote}
Before the new one: say the Mandarin for lucky by a fluke, then the Mandarin for a misfortune.
\end{quote}

\subsection*{You'll want to know: shīshǒu}
\emph{shīshǒu} — \textquotedblleft{}to lose one's grip; to do by a slip of the hand\textquotedblright{}.

\begin{quote}
\emph{yàoshi nǐ shīshǒu dǎpò le bēizi jiù zài mǎi yí ge} — If your hand slips and you break a cup, buy another one.
\end{quote}

Its characters come later; for now it is a word for the ear.

\begin{cousinweb}
\emph{shīshǒu} is built from \emph{shī} 'to lose' and \zh{手} 'the hand', both inherited from Middle Chinese, \emph{syit} and \emph{syuwX} (Guangyun, in Baxter's transcription), which glosses the first 'a mistake' and 'to let go': the hand lets it go.
\end{cousinweb}

\begin{grammarlens}[title={What you've built}]
One more word for saying what would happen.
\end{grammarlens}

\subsection*{Your turn}
\begin{itemize}
\raggedright
  \item \emph{Say it:} \emph{shīshǒu}
  \item \emph{Say it:} \emph{shīshǒu}, once more
  \item \emph{Say it:} \emph{yàoshi nǐ shīshǒu dǎpò le bēizi jiù zài mǎi yí ge} — If your hand slips and you break a cup, buy another one.
  \item \emph{Recall:} say the Mandarin for lucky by a fluke, then the Mandarin for a misfortune, then say \emph{shīshǒu} again
\end{itemize}

\begin{tcolorbox}[breakable,colback=teal!4,colframe=teal!35!black,title={Before you move on}]
What does shīshǒu mean? (\textquotedblleft{}To lose one's grip; to do by a slip of the hand\textquotedblright{}.) Say it once more.
\end{tcolorbox}
\section[(dialogue)]{Looking back: from How Heavy? to A Lucky Escape}

\label{lesson:ZH-R531-review-526-531}



\begin{quote}
Say the Mandarin for to weigh on a scale and a kilogram.
\end{quote}

\subsection*{Your turn}
\begin{itemize}
\raggedright
  \item \emph{Say it:} to weigh on a scale, a kilogram, how heavy something is, in large amounts, a small amount of
  \item \emph{Say it:} to place an order, to wrap goods for sale, a member of a club, a favourable deal, a checked pattern
  \item \emph{Say it:} as luck would have it, to be obliged to, to draw lots, to make a bet, happening by chance
  \item \emph{Say it:} to be out of luck, a person's luck, to be in luck, fortunate, fate
  \item \emph{Say it:} a lottery ticket, to draw for prizes, to win in a draw, a prize you win, the jackpot
  \item \emph{Say it:} to get out of danger, all but, lucky by a fluke, a misfortune, to lose one's grip
  \item \emph{Read it:} \zh{量} \zh{运} — each one aloud, then a word that holds it
\end{itemize}

\begin{tcolorbox}[breakable,colback=teal!4,colframe=teal!35!black,title={Before you move on}]
To weigh on a scale? (\emph{chēng}.) A kilogram? (\emph{gōngjīn}.) How heavy something is? (\textbf{\zh{重量}}, \emph{zhòngliàng}.) In large amounts? (\textbf{\zh{大量}}, \emph{dàliàng}.) A small amount of? (\textbf{\zh{少量}}, \emph{shǎoliàng}.) To place an order? (\emph{xiàdān}.) To wrap goods for sale? (\emph{bāozhuāng}.) A member of a club? (\emph{huìyuán}.) A favourable deal? (\emph{yōuhuì}.) A checked pattern? (\textbf{\zh{格子}}, \emph{gézi}.) As luck would have it? (\emph{pèngqiǎo}.) To be obliged to? (\textbf{\zh{不得不}}, \emph{bùdébù}.) To draw lots? (\emph{chōuqiān}.) To make a bet? (\emph{dǎdǔ}.) Happening by chance? (\emph{ǒurán}.) To be out of luck? (\emph{dǎoméi}.) A person's luck? (\textbf{\zh{运气}}, \emph{yùnqi}.) To be in luck? (\textbf{\zh{走运}}, \emph{zǒuyùn}.) Fortunate? (\emph{xìngyùn}.) Fate? (\emph{mìngyùn}.) A lottery ticket? (\emph{cǎipiào}.) To draw for prizes? (\emph{chōujiǎng}.) To win in a draw? (\emph{zhòngjiǎng}.) A prize you win? (\emph{jiǎngpǐn}.) The jackpot? (\emph{tóujiǎng}.) To get out of danger? (\emph{tuōxiǎn}.) All but? (\emph{jīhū}.) Lucky by a fluke? (\emph{jiǎoxìng}.) A misfortune? (\emph{búxìng}.) To lose one's grip? (\emph{shīshǒu}.)
\end{tcolorbox}
